% ============================================================
%  Semaine 2 — HTML5 : structure & sémantique  [Module B]
%  PRÉSENTATION (Beamer) — cfpt·i | Anne Terrier
%  Aligné sur PlanCours_Pass_Web_final et sur S02_polycop
%  (les formulaires et l'accessibilité sont traités en semaine 3)
%
%  Images attendues dans ./images/ : aucune pour l'instant.
%  Une image absente est remplacée par un cadre de réservation.
% ============================================================
\documentclass[aspectratio=169,11pt]{beamer}

\usepackage[T1]{fontenc}
\usepackage[utf8]{inputenc}
\usepackage[french]{babel}
\usepackage[scaled=0.92]{helvet}
\renewcommand{\familydefault}{\sfdefault}

\usepackage{graphicx}
\usepackage{tikz}
\usetikzlibrary{positioning, arrows.meta, decorations.pathreplacing}
\usepackage{booktabs}
\usepackage{tabularx}
\usepackage{colortbl}
\usepackage{listings}

% ---------------------------------------------------------------
% Correction de l'alignement des en-tetes de tableau
%   1) \aboverulesep / \belowrulesep a 0 : supprime le lisere blanc
%      entre le filet booktabs et la bande coloree du \rowcolor.
%   2) \hdr : force l'alignement de la cellule d'en-tete. Sans cela,
%      une colonne p{}, m{} ou X place le texte trop bas dans la bande.
% ---------------------------------------------------------------
\setlength{\aboverulesep}{0pt}
\setlength{\belowrulesep}{0pt}
\newcommand{\hdr}[1]{\multicolumn{1}{l}{\color{white}\bfseries #1}}
\newcommand{\hdrc}[1]{\multicolumn{1}{c}{\color{white}\bfseries #1}}
\newcommand{\hdrr}[1]{\multicolumn{1}{r}{\color{white}\bfseries #1}}

% ---------------------------------------------------------------
% Palette cfpt·i
% ---------------------------------------------------------------
\definecolor{cfptYellow}{HTML}{F5C400}
\definecolor{cfptBlack}{HTML}{1A1A1A}
\definecolor{cfptDarkGrey}{HTML}{555555}
\definecolor{cfptGrey}{HTML}{F2F2F2}
\definecolor{accentBlue}{HTML}{1D4E89}
\definecolor{lightBlue}{HTML}{D6E4F0}
\definecolor{codeGrey}{HTML}{F5F5F5}
\definecolor{successGreen}{HTML}{1E7145}
\definecolor{warningOrange}{HTML}{C55A11}
\definecolor{alertRed}{HTML}{A32020}

% Alias (repris du polycopié)
\colorlet{uniInk}{cfptBlack}
\colorlet{uniNavy}{accentBlue}
\colorlet{uniTeal}{successGreen}
\colorlet{uniGold}{warningOrange}
\colorlet{uniGrey}{cfptDarkGrey}
\colorlet{uniPaper}{cfptGrey}
\colorlet{uniRule}{cfptDarkGrey!35}

% ---------------------------------------------------------------
% Image robuste : cadre de réservation si le fichier est absent
% ---------------------------------------------------------------
\newcommand{\imageslide}[2][0.80\textwidth]{%
  \IfFileExists{#2}%
    {\includegraphics[width=#1]{#2}}%
    {\setlength{\fboxrule}{0.6pt}%
     \textcolor{cfptDarkGrey!60}{\fbox{\parbox[c][3.6cm][c]{#1}{%
       \centering\small\color{cfptDarkGrey}%
       [ image à insérer ]\\[4pt]\texttt{\detokenize{#2}}}}}}%
}

% ---------------------------------------------------------------
% Thème Beamer cfpt·i
% ---------------------------------------------------------------
\usetheme{default}
\usecolortheme{default}
\setbeamercolor{palette primary}{bg=accentBlue, fg=white}
\setbeamercolor{structure}{fg=accentBlue}
\setbeamercolor{frametitle}{bg=accentBlue, fg=white}
\setbeamercolor{title}{fg=white}
\setbeamercolor{block title}{bg=accentBlue, fg=white}
\setbeamercolor{block body}{bg=lightBlue!30, fg=cfptBlack}
\setbeamercolor{block title alerted}{bg=cfptYellow!85!black, fg=cfptBlack}
\setbeamercolor{block body alerted}{bg=cfptYellow!20, fg=cfptBlack}
\setbeamercolor{block title example}{bg=successGreen, fg=white}
\setbeamercolor{block body example}{bg=successGreen!8, fg=cfptBlack}
\setbeamercolor{itemize item}{fg=cfptYellow!85!black}
\setbeamercolor{itemize subitem}{fg=cfptDarkGrey}
\setbeamercolor{enumerate item}{fg=accentBlue}
\setbeamercolor{normal text}{fg=cfptBlack}
\setbeamercolor{footline}{fg=white}

\setbeamertemplate{itemize item}{\textbullet}
\setbeamertemplate{itemize subitem}{--}
\setbeamertemplate{navigation symbols}{}
\setbeamertemplate{blocks}[rounded][shadow=false]

% Titre de frame : bandeau pleine largeur
\setbeamertemplate{frametitle}{%
  \nointerlineskip%
  \begin{beamercolorbox}[wd=\paperwidth, ht=2.6ex, dp=1.2ex, leftskip=0.3cm]{frametitle}%
    \usebeamerfont{frametitle}\bfseries\insertframetitle%
  \end{beamercolorbox}%
}

% Pied de page
\setbeamertemplate{footline}{%
  \leavevmode%
  \hbox{%
    \begin{beamercolorbox}[wd=.32\paperwidth, ht=2.4ex, dp=1ex, leftskip=0.3cm]{palette primary}%
      \usebeamerfont{author in head/foot}\tiny Anne Terrier%
    \end{beamercolorbox}%
    \begin{beamercolorbox}[wd=.44\paperwidth, ht=2.4ex, dp=1ex, center]{palette primary}%
      \usebeamerfont{title in head/foot}\tiny Semaine 2 --- HTML5 : structure \& sémantique%
    \end{beamercolorbox}%
    \begin{beamercolorbox}[wd=.24\paperwidth, ht=2.4ex, dp=1ex, right, rightskip=0.3cm]{palette primary}%
      \usebeamerfont{date in head/foot}\tiny\insertframenumber\,/\,\inserttotalframenumber%
    \end{beamercolorbox}%
  }%
}

% Diapositive de séparation de partie
\newcommand{\partieslide}[1]{%
{%
\setbeamertemplate{footline}{}%
\begin{frame}[plain]
  \begin{tikzpicture}[remember picture, overlay]
    \fill[accentBlue] (current page.north west) rectangle (current page.south east);
    \fill[cfptYellow] ([yshift=-0.15cm]current page.center -| current page.west)
      rectangle ([yshift=-0.3cm]current page.center -| current page.east);
    \node[text=white, font=\Huge\bfseries] at ([yshift=0.9cm]current page.center) {#1};
  \end{tikzpicture}
\end{frame}
}}

% Diapositive « à vous de jouer » (fond vert)
\newcommand{\exoslide}[1]{%
{%
\setbeamertemplate{footline}{}%
\begin{frame}[plain]
  \begin{tikzpicture}[remember picture, overlay]
    \fill[successGreen] (current page.north west) rectangle (current page.south east);
    \fill[cfptYellow] ([yshift=-0.15cm]current page.center -| current page.west)
      rectangle ([yshift=-0.3cm]current page.center -| current page.east);
    \node[text=white, font=\Huge\bfseries] at ([yshift=1.3cm]current page.center)
      {À vous de jouer !};
    \node[text=white, font=\large, text width=0.8\paperwidth, align=center]
      at ([yshift=-0.9cm]current page.center) {#1};
  \end{tikzpicture}
\end{frame}
}}

% ---------------------------------------------------------------
% Listings
% ---------------------------------------------------------------
\lstdefinelanguage{HTML5}{
  keywords={<!DOCTYPE, html, head, body, header, nav, main, section, article,
    aside, footer, h1, h2, h3, h4, h5, h6, p, a, img, figure, figcaption,
    ul, ol, li, table, caption, thead, tbody, tr, th, td, div, span, meta,
    title, link, script, strong, em, br, hr, blockquote},
  sensitive=false,
  morecomment=[s]{<!--}{-->},
  morestring=[b]",
  keywordstyle=\color{accentBlue}\bfseries,
  stringstyle=\color{warningOrange},
  commentstyle=\color{cfptDarkGrey}\itshape,
}

\lstset{
  basicstyle=\ttfamily\scriptsize,
  backgroundcolor=\color{codeGrey},
  frame=single, framerule=0.4pt,
  rulecolor=\color{cfptDarkGrey!40},
  breaklines=true,
  keywordstyle=\color{accentBlue}\bfseries,
  commentstyle=\color{cfptDarkGrey}\itshape,
  stringstyle=\color{successGreen},
  showstringspaces=false,
  extendedchars=true,
  literate=%
    {é}{{\'e}}1 {è}{{\`e}}1 {ê}{{\^e}}1 {ë}{{\"e}}1
    {É}{{\'E}}1 {È}{{\`E}}1 {Ê}{{\^E}}1
    {à}{{\`a}}1 {â}{{\^a}}1 {À}{{\`A}}1
    {î}{{\^i}}1 {ï}{{\"i}}1 {ô}{{\^o}}1
    {û}{{\^u}}1 {ù}{{\`u}}1 {ü}{{\"u}}1
    {ç}{{\c c}}1 {Ç}{{\c C}}1 {œ}{{\oe}}1
}

% ---------------------------------------------------------------
% Métadonnées
% ---------------------------------------------------------------
\title{HTML5 --- Structure \& sémantique}
\subtitle{Développement web --- Classe passerelle}
\author{Anne Terrier}
\institute{cfpt\textperiodcentered i --- École d'informatique}
\date{Semaine 2}

% ===============================================================
\begin{document}
% ===============================================================

% --- Diapo de titre ---
{
\setbeamertemplate{footline}{}
\begin{frame}[plain]
  \begin{tikzpicture}[remember picture, overlay]
    \fill[accentBlue] (current page.north west) rectangle ([yshift=-3.2cm]current page.north east);
    \node[anchor=north west, text=white, font=\huge\bfseries, text width=0.9\paperwidth]
      at ([shift={(0.8cm,-0.8cm)}]current page.north west)
      {HTML5 --- Structure \& sémantique};
    \node[anchor=north west, text=lightBlue, font=\large]
      at ([shift={(0.85cm,-2.55cm)}]current page.north west)
      {Développement web --- Classe passerelle};

    % accent jaune
    \fill[cfptYellow] ([yshift=-3.2cm]current page.north west) rectangle ([yshift=-3.35cm]current page.north east);

    \node[anchor=south west, font=\large\bfseries, text=accentBlue]
      at ([shift={(0.8cm,1.8cm)}]current page.south west) {Semaine 2 \quad\textbar\quad Module B};
    \node[anchor=south west, font=\normalsize, text=cfptDarkGrey]
      at ([shift={(0.8cm,1.1cm)}]current page.south west)
      {Anne Terrier \quad\textbar\quad cfpt\textperiodcentered i --- École d'informatique};
  \end{tikzpicture}
\end{frame}
}

% --- Plan ---
\begin{frame}{Au programme aujourd'hui}
  \begin{columns}[T]
    \begin{column}{0.48\textwidth}
      \begin{enumerate}
        \item Le HTML, langage de balisage
        \item La structure d'un document
        \item Les éléments sémantiques
        \item Structurer le texte
      \end{enumerate}
      \vspace{0.25cm}
      \begin{exampleblock}{Au milieu}
        Un exercice : \textbf{votre première page web}.
      \end{exampleblock}
    \end{column}
    \begin{column}{0.48\textwidth}
      \begin{enumerate}
        \setcounter{enumi}{4}
        \item Liens \& images
        \item Listes \& tableaux
        \item Valider son HTML
      \end{enumerate}
      \vspace{0.25cm}
      \begin{block}{Objectif de la semaine}
        Écrire une page HTML5 valide et sémantique, et savoir corriger les erreurs signalées par le validateur.
      \end{block}
    \end{column}
  \end{columns}
\end{frame}

\begin{frame}{Où en sommes-nous ?}
  \begin{block}{La semaine dernière}
    Le modèle client / serveur, les trois langages du navigateur, l'environnement de travail et le premier dépôt Git.
  \end{block}
  \vspace{0.3cm}
  \begin{center}
  \begin{tikzpicture}[font=\small,
    b/.style={draw, rounded corners=2pt, minimum width=3.2cm,
              minimum height=1.15cm, align=center}]
    \node[b, draw=warningOrange, fill=warningOrange!12] (h) at (0,0)
      {\textbf{HTML}\\{\footnotesize la structure}};
    \node[b, draw=accentBlue, fill=accentBlue!8, right=0.6cm of h] (c)
      {\textbf{CSS}\\{\footnotesize la présentation}};
    \node[b, draw=successGreen, fill=successGreen!8, right=0.6cm of c] (j)
      {\textbf{JavaScript}\\{\footnotesize le comportement}};
    \node[below=0.12cm of h, font=\footnotesize\bfseries, text=warningOrange]
      {semaines 2 et 3};
    \node[below=0.12cm of c, font=\footnotesize, text=cfptDarkGrey]
      {semaines 4 et 5};
    \node[below=0.12cm of j, font=\footnotesize, text=cfptDarkGrey]
      {semaines 6 et 7};
  \end{tikzpicture}
  \end{center}
  \vspace{0.35cm}
  \begin{alertblock}{Aujourd'hui, on ne fait que du HTML}
    Votre page sera laide. C'est normal, et c'est même voulu : tant qu'il n'y a pas de style, on ne peut juger que la \textbf{structure}.
  \end{alertblock}
\end{frame}

% ===============================================================
\section{Le HTML, langage de balisage}
% ===============================================================
\partieslide{1. Le HTML, langage de balisage}

\begin{frame}{Un langage de balisage}
  \begin{block}{Définition}
    Le \textbf{HTML} (\textit{HyperText Markup Language}) décrit la \textbf{structure} et le \textbf{contenu} d'une page web. Une page est un simple fichier texte d'extension \texttt{.html}.
  \end{block}
  \vspace{0.3cm}
  \begin{alertblock}{Ce n'est pas un langage de programmation}
    Ni variables, ni conditions, ni boucles : on \textbf{balise} le contenu pour dire au navigateur ce que chaque morceau représente.
    \vspace{0.15cm}

    Les notions de programmation arriveront avec JavaScript, en \textbf{semaine 6}.
  \end{alertblock}
\end{frame}

\begin{frame}{Anatomie d'une balise}
  \vspace{0.4cm}
  \begin{center}
  \begin{tikzpicture}[font=\small]
    \node[font=\ttfamily\Large] (code) {%
      \textcolor{accentBlue}{<p>}%
      \textcolor{cfptBlack}{~Bonjour tout le monde~}%
      \textcolor{accentBlue}{</p>}};
    \node[below=1.1cm of code, xshift=-3.4cm, text=accentBlue, font=\footnotesize\bfseries] (lo) {Balise ouvrante};
    \node[below=1.1cm of code, xshift=0cm, text=cfptDarkGrey, font=\footnotesize\bfseries] (ct) {Contenu};
    \node[below=1.1cm of code, xshift=3.3cm, text=accentBlue, font=\footnotesize\bfseries] (lf) {Balise fermante};
    \draw[-{Stealth[length=5pt]}, accentBlue!70, thick] (lo.north) -- ++(0,0.65);
    \draw[-{Stealth[length=5pt]}, cfptDarkGrey!70, thick] (ct.north) -- ++(0,0.65);
    \draw[-{Stealth[length=5pt]}, accentBlue!70, thick] (lf.north) -- ++(0,0.65);
  \end{tikzpicture}
  \end{center}
  \vspace{0.3cm}
  \begin{columns}[T]
    \begin{column}{0.49\textwidth}
      \begin{block}{Un élément}
        L'ensemble ouvrante + contenu + fermante. La fermante se reconnaît à sa barre oblique.
      \end{block}
    \end{column}
    \begin{column}{0.49\textwidth}
      \begin{exampleblock}{Balises auto-fermantes}
        Sans contenu, elles s'écrivent seules : \texttt{<img>}, \texttt{<br>}, \texttt{<meta>}, \texttt{<link>}.
      \end{exampleblock}
    \end{column}
  \end{columns}
\end{frame}

\begin{frame}[fragile]{Les attributs}
  Un attribut précise le comportement d'une balise. Il s'écrit dans la balise \textbf{ouvrante}, sous la forme \texttt{nom="valeur"} :
  \vspace{0.3cm}
\begin{lstlisting}[language=HTML5]
<a href="https://www.hepia.ch" target="_blank">Site de la hepia</a>
\end{lstlisting}
  \vspace{0.3cm}
  \begin{itemize}
    \item \texttt{href} indique la \textbf{destination} du lien
    \item \texttt{target} demande l'ouverture dans un \textbf{nouvel onglet}
  \end{itemize}
  \vspace{0.3cm}
  \begin{center}
    \small Une balise \textbf{fermante} ne porte jamais d'attribut.
  \end{center}
\end{frame}

\begin{frame}{L'imbrication : un document est un arbre}
  \vspace{0.15cm}
  \begin{center}
  \begin{tikzpicture}[font=\scriptsize,
    n/.style={draw=uniNavy, fill=uniNavy!8, rounded corners=2pt,
              inner sep=3.5pt, font=\ttfamily\scriptsize},
    f/.style={draw=uniGrey!55, fill=uniPaper, rounded corners=2pt,
              inner sep=3.5pt, font=\ttfamily\scriptsize},
    e/.style={uniGrey!70, thick}]
    \node[n] (html) at (0,0) {<html>};
    \node[n] (head) at (-2.6,-1.0) {<head>};
    \node[n] (body) at (2.1,-1.0) {<body>};
    \node[f] (title) at (-2.6,-2.0) {<title>};
    \node[n] (h1)   at (0.4,-2.0) {<h1>};
    \node[n] (ul)   at (2.8,-2.0) {<ul>};
    \node[f] (li1)  at (1.9,-3.0) {<li>};
    \node[f] (li2)  at (3.7,-3.0) {<li>};
    \draw[e] (html) -- (head);  \draw[e] (html) -- (body);
    \draw[e] (head) -- (title); \draw[e] (body) -- (h1);
    \draw[e] (body) -- (ul);    \draw[e] (ul) -- (li1);
    \draw[e] (ul) -- (li2);
  \end{tikzpicture}
  \end{center}
  \vspace{0.15cm}
  \begin{block}{Pourquoi ça compte}
    Cet arbre sera parcouru par le navigateur, par les moteurs de recherche, et --- dès la semaine 7 --- par votre code JavaScript.
  \end{block}
\end{frame}

\begin{frame}{La règle d'imbrication}
  \begin{alertblock}{Premier ouvert, dernier fermé}
    Les balises se ferment dans l'ordre inverse de leur ouverture, comme des parenthèses. Elles ne peuvent pas se chevaucher.
  \end{alertblock}
  \vspace{0.4cm}
  \begin{center}
    \large
    \textcolor{successGreen}{\textbf{OUI}}\quad
    \texttt{<p>Un <strong>mot</strong> important</p>}

    \vspace{0.3cm}
    \textcolor{alertRed}{\textbf{NON}}\quad
    \texttt{<p>Un <strong>mot</p></strong> important}
  \end{center}
  \vspace{0.4cm}
  \begin{exampleblock}{Le réflexe qui évite tout}
    \textbf{Indenter} : un niveau d'indentation par niveau d'imbrication. Une balise oubliée devient alors visible à l'\oe il nu.
  \end{exampleblock}
\end{frame}

\begin{frame}[fragile]{Le HTML ignore votre mise en forme du texte}
  \begin{alertblock}{Espaces et retours à la ligne}
    Plusieurs espaces, les tabulations et les retours à la ligne sont \textbf{réduits à un seul espace} à l'affichage.
    \vspace{0.15cm}

    Appuyer sur Entrée ne crée pas de paragraphe : c'est un nouveau \texttt{<p>} qui le fait.
  \end{alertblock}
  \vspace{0.2cm}
\begin{lstlisting}[language=HTML5]
<!-- Un commentaire : invisible dans la page,
     mais lisible par tous via Ctrl+U -->
<p>Ce      texte    s'affichera    normalement.</p>
\end{lstlisting}
  \vspace{0.15cm}
  \begin{center}
    \small Un commentaire sert à \textbf{expliquer} son code, jamais à y cacher une information.
  \end{center}
\end{frame}

% ===============================================================
\section{La structure d'un document HTML5}
% ===============================================================
\partieslide{2. La structure d'un document}

\begin{frame}[fragile]{Le squelette d'une page HTML5}
\begin{lstlisting}[language=HTML5]
<!DOCTYPE html>
<html lang="fr">
<head>
  <meta charset="UTF-8">
  <meta name="viewport" content="width=device-width, initial-scale=1.0">
  <title>MiniForum</title>
  <link rel="stylesheet" href="style.css">
</head>
<body>
  <!-- Tout le contenu visible de la page -->
</body>
</html>
\end{lstlisting}
  \vspace{0.15cm}
  \begin{center}
    \small À connaître par c\oe ur --- il sera demandé au contrôle écrit de la semaine 12.
  \end{center}
\end{frame}

\begin{frame}{Ce que contient le \texttt{<head>}}
  \begin{center}
  \small
  \begin{tabularx}{0.97\textwidth}{p{3.9cm}X}
    \toprule
    \rowcolor{accentBlue}
    \hdr{Balise} & \hdr{Ce qu'elle fait} \\
    \midrule
    \rowcolor{cfptGrey}
    \texttt{<meta charset="UTF-8">} & L'encodage. Sans elle, les accents deviennent illisibles \\
    \texttt{<meta name="viewport">} & Affichage correct sur téléphone (responsive, semaine 5) \\
    \rowcolor{cfptGrey}
    \texttt{<title>} & Le texte de l'onglet, des favoris et des résultats de recherche \\
    \texttt{<link rel="stylesheet">} & Rattache la feuille de style CSS \\
    \bottomrule
  \end{tabularx}
  \end{center}
  \vspace{0.3cm}
  \begin{alertblock}{\texttt{<head>} et \texttt{<body>} : ne pas confondre}
    Le \texttt{<head>} contient des informations \emph{sur} la page. Le \texttt{<body>} contient tout le \textbf{visible}. Un titre écrit dans le \texttt{<head>} n'apparaîtra jamais dans la page.
  \end{alertblock}
\end{frame}

% ===============================================================
\section{Les éléments sémantiques}
% ===============================================================
\partieslide{3. Les éléments sémantiques}

\begin{frame}{Décrire le rôle de chaque zone}
  \vspace{0.1cm}
  \begin{center}
  \begin{tikzpicture}[font=\scriptsize, every node/.style={align=center}, scale=0.92, transform shape]
    \draw[draw=accentBlue, fill=accentBlue!12, rounded corners=2pt] (0,0) rectangle (10,-0.75);
    \node[text=accentBlue] at (5,-0.375) {\texttt{\textbf{<header>}} --- en-tête : titre du site, logo};

    \draw[draw=warningOrange, fill=warningOrange!12, rounded corners=2pt] (0,-0.9) rectangle (10,-1.55);
    \node[text=warningOrange] at (5,-1.22) {\texttt{\textbf{<nav>}} --- menu de navigation};

    \draw[draw=successGreen, fill=successGreen!8, rounded corners=2pt] (0,-1.7) rectangle (7.4,-4.5);
    \node[text=successGreen, font=\scriptsize\bfseries] at (1.0,-1.95) {\texttt{<main>}};
    \draw[draw=successGreen!70, fill=white, rounded corners=2pt] (0.3,-2.2) rectangle (7.1,-3.25);
    \node at (3.7,-2.72) {\texttt{\textbf{<article>}} --- un sujet du forum};
    \draw[draw=successGreen!70, fill=white, rounded corners=2pt] (0.3,-3.4) rectangle (7.1,-4.3);
    \node at (3.7,-3.85) {\texttt{\textbf{<article>}} --- un autre sujet};

    \draw[draw=cfptDarkGrey, fill=cfptGrey, rounded corners=2pt] (7.6,-1.7) rectangle (10,-4.5);
    \node[text=cfptDarkGrey] at (8.8,-3.1) {\texttt{\textbf{<aside>}}\\[2pt]contenu\\annexe};

    \draw[draw=accentBlue, fill=accentBlue!12, rounded corners=2pt] (0,-4.65) rectangle (10,-5.3);
    \node[text=accentBlue] at (5,-4.97) {\texttt{\textbf{<footer>}} --- pied de page : mentions, contact};
  \end{tikzpicture}
  \end{center}
\end{frame}

\begin{frame}{Pourquoi la sémantique ?}
  Avant HTML5, les pages n'étaient qu'un empilement de \texttt{<div>}. Les balises sémantiques disent ce que le contenu \textbf{est}, pas comment il s'affiche.
  \vspace{0.3cm}
  \begin{itemize}
    \item les \textbf{humains} qui relisent le code --- vous, dans trois semaines
    \item les \textbf{moteurs de recherche}, qui repèrent le contenu principal
    \item les \textbf{lecteurs d'écran}, qui permettent de sauter directement au contenu
  \end{itemize}
  \vspace{0.35cm}
  \begin{alertblock}{Déroutant, mais essentiel}
    Remplacer un \texttt{<div>} par un \texttt{<main>} \textbf{ne change rien à l'affichage}. Le bénéfice est invisible --- et bien réel. C'est aussi pourquoi le TP se fait sans CSS : rien ne doit masquer la structure.
  \end{alertblock}
\end{frame}

\begin{frame}{\texttt{<section>} ou \texttt{<article>} ?}
  \begin{block}{Le test}
    Ce contenu garderait-il du sens \textbf{sorti de la page}, publié seul ailleurs ?
  \end{block}
  \vspace{0.4cm}
  \begin{columns}[T]
    \begin{column}{0.49\textwidth}
      \begin{exampleblock}{Oui $\rightarrow$ \texttt{<article>}}
        Un message de forum, un billet de blog, une fiche produit. Autonome.
      \end{exampleblock}
    \end{column}
    \begin{column}{0.49\textwidth}
      \begin{exampleblock}{Non $\rightarrow$ \texttt{<section>}}
        « Sujets récents » n'a de sens que dans le forum qui la contient.
      \end{exampleblock}
    \end{column}
  \end{columns}
  \vspace{0.4cm}
  \begin{center}
    \small Et le \texttt{<div>} ? Il reste utile pour regrouper des éléments \textbf{uniquement} en vue de la mise en forme, quand il n'y a aucun sens à exprimer.
  \end{center}
\end{frame}

% ===============================================================
\section{Structurer le texte}
% ===============================================================
\partieslide{4. Structurer le texte}

\begin{frame}[fragile]{Titres et paragraphes}
\begin{lstlisting}[language=HTML5]
<h1>MiniForum</h1>
<h2>Sujets récents</h2>
<h3>Bienvenue sur le forum</h3>

<p>Un paragraphe avec du <strong>contenu important</strong>
   et un mot en <em>emphase</em>.</p>
\end{lstlisting}
  \vspace{0.2cm}
  \begin{alertblock}{\texttt{<strong>} plutôt que \texttt{<b>}, \texttt{<em>} plutôt que \texttt{<i>}}
    \texttt{<b>} et \texttt{<i>} décrivent une \textit{apparence} ; \texttt{<strong>} et \texttt{<em>} décrivent un \textit{sens}. Le rendu est identique --- mais un lecteur d'écran change d'intonation sur \texttt{<strong>} et ignore \texttt{<b>}.
  \end{alertblock}
\end{frame}

\begin{frame}{La hiérarchie des titres}
  Les titres construisent le \textbf{plan} du document, comme dans un rapport. Ce n'est pas un outil de mise en forme.
  \vspace{0.25cm}
  \begin{columns}[T]
    \begin{column}{0.49\textwidth}
      \begin{exampleblock}{Plan correct}
        \texttt{h1} MiniForum\\
        \quad\texttt{h2} Sujets récents\\
        \quad\quad\texttt{h3} Bienvenue\\
        \quad\quad\texttt{h3} Installer Node.js\\
        \quad\texttt{h2} Nouveau sujet
      \end{exampleblock}
    \end{column}
    \begin{column}{0.49\textwidth}
      \begin{alertblock}{Plan incorrect}
        \texttt{h1} MiniForum\\
        \quad\texttt{h3} Sujets récents \textcolor{alertRed}{\textbf{(saut)}}\\
        \quad\quad\texttt{h2} Bienvenue\\
        \texttt{h1} Nouveau sujet \textcolor{alertRed}{\textbf{(2\textsuperscript{e} h1)}}
      \end{alertblock}
    \end{column}
  \end{columns}
  \vspace{0.3cm}
  \begin{center}
    \small Deux règles : \textbf{un seul \texttt{<h1>}} par page, et \textbf{jamais de saut de niveau}.
    Choisir un \texttt{<h3>} « parce qu'il fait la bonne taille » est une erreur --- la taille, c'est du CSS.
  \end{center}
\end{frame}

\begin{frame}{Éléments en bloc et éléments en ligne}
  \begin{columns}[T]
    \begin{column}{0.49\textwidth}
      \begin{block}{En bloc}
        Occupe toute la largeur et provoque un retour à la ligne.
        \vspace{0.15cm}

        \texttt{<p>}, \texttt{<h1>}, \texttt{<ul>}, \texttt{<section>}, \texttt{<div>}
      \end{block}
    \end{column}
    \begin{column}{0.49\textwidth}
      \begin{exampleblock}{En ligne}
        Reste dans le fil du texte, sans le couper.
        \vspace{0.15cm}

        \texttt{<a>}, \texttt{<strong>}, \texttt{<em>}, \texttt{<img>}, \texttt{<span>}
      \end{exampleblock}
    \end{column}
  \end{columns}
  \vspace{0.5cm}
  \begin{center}
    \small Cela explique pourquoi deux paragraphes se suivent verticalement\\
    alors que deux liens restent côte à côte.
    \vspace{0.2cm}

    \small\itshape Cette distinction deviendra centrale en semaine 4, avec le modèle de boîte du CSS.
  \end{center}
\end{frame}

% ===============================================================
\section{Liens \& images}
% ===============================================================
\partieslide{5. Liens \& images}

\begin{frame}[fragile]{Le lien hypertexte}
\begin{lstlisting}[language=HTML5]
<!-- Lien externe : URL complète -->
<a href="https://www.hepia.ch">Site de la hepia</a>

<!-- Vers un autre fichier du site -->
<a href="sujet.html">Voir le sujet</a>

<!-- Vers un élément de la page courante -->
<a href="#nouveau-sujet">Créer un sujet</a>
\end{lstlisting}
  \vspace{0.2cm}
  \begin{alertblock}{Écrire un bon intitulé}
    Le texte du lien doit décrire sa destination \textbf{hors contexte} : les lecteurs d'écran savent lister tous les liens d'une page, et une série de « cliquez ici » n'informe personne.
    \vspace{0.15cm}

    \textcolor{successGreen}{\textbf{OUI}} \texttt{Lire le sujet complet} \qquad
    \textcolor{alertRed}{\textbf{NON}} \texttt{cliquez ici}
  \end{alertblock}
\end{frame}

\begin{frame}[fragile]{Chemins relatifs : la source d'erreur n\textsuperscript{o}\,1}
  \begin{columns}[T]
    \begin{column}{0.42\textwidth}
\begin{lstlisting}
miniforum/
  |-- index.html
  |-- sujet.html
  |-- style.css
  \-- images/
       \-- logo.png
\end{lstlisting}
    \end{column}
    \begin{column}{0.56\textwidth}
      \small
      \begin{tabularx}{\linewidth}{p{2.5cm}X}
        \toprule
        \rowcolor{accentBlue}
        \hdr{Écriture} & \hdr{Signification} \\
        \midrule
        \rowcolor{cfptGrey}
        \texttt{sujet.html} & même dossier \\
        \texttt{images/logo.png} & sous-dossier \\
        \rowcolor{cfptGrey}
        \texttt{../style.css} & dossier parent \\
        \texttt{https://\dots} & autre serveur \\
        \bottomrule
      \end{tabularx}
    \end{column}
  \end{columns}
  \vspace{0.3cm}
  \begin{alertblock}{Trois pièges}
    \textbf{La casse compte} sur un serveur Linux (\texttt{Logo.png} $\neq$ \texttt{logo.png}) : ça marche chez vous, ça casse en ligne. \quad
    \textbf{Pas d'espaces ni d'accents} dans les noms de fichiers. \quad
    \textbf{Jamais} de chemin \texttt{C:/Users/\dots}
  \end{alertblock}
\end{frame}

\begin{frame}[fragile]{L'attribut \texttt{id}}
  Un \texttt{id} donne un \textbf{nom unique} à un élément de la page :
\begin{lstlisting}[language=HTML5]
<a href="#nouveau-sujet">Créer un sujet</a>

<section id="nouveau-sujet"> ... </section>
\end{lstlisting}
  \vspace{0.15cm}
  \begin{columns}[T]
    \begin{column}{0.49\textwidth}
      \begin{alertblock}{Règle d'or}
        Un \texttt{id} est \textbf{unique} dans la page. Le validateur signale les doublons comme des erreurs.
      \end{alertblock}
    \end{column}
    \begin{column}{0.49\textwidth}
      \begin{block}{Un point d'accroche pour l'année}
        Le \texttt{<label>} (sem. 3), le \textbf{CSS} (sem. 4), puis \textbf{JavaScript} (sem. 7) s'en serviront tous.
      \end{block}
    \end{column}
  \end{columns}
\end{frame}

\begin{frame}[fragile]{Les images}
\begin{lstlisting}[language=HTML5]
<img src="images/logo.png" alt="Logo du MiniForum">

<figure>
  <img src="images/schema.png" alt="Organisation des pages du forum">
  <figcaption>Organisation des pages du MiniForum</figcaption>
</figure>
\end{lstlisting}
  \vspace{0.2cm}
  \begin{alertblock}{L'attribut \texttt{alt} n'est pas optionnel}
    Il est lu à voix haute par les lecteurs d'écran, et s'affiche si l'image ne charge pas. Le validateur signale son absence comme une \textbf{erreur}.
    \vspace{0.15cm}

    \textcolor{successGreen}{\textbf{OUI}} \texttt{alt="Logo du MiniForum"} \qquad
    \textcolor{alertRed}{\textbf{NON}} \texttt{alt="image"}
  \end{alertblock}
\end{frame}

% ===============================================================
% EXERCICE INTERMÉDIAIRE
% ===============================================================
\exoslide{Exercice intermédiaire --- votre première page web}

\begin{frame}{Exercice --- votre première page web}
  \begin{exampleblock}{Objectif}
    Écrire seul·e une page HTML5 valide et sémantique. Pas de MiniForum ici : une courte page de présentation personnelle.
  \end{exampleblock}
  \begin{block}{Consignes --- fichier \texttt{ma-page.html}}
    \begin{enumerate}
      \item Le squelette écrit \textbf{à la main} : \texttt{<!DOCTYPE>}, \texttt{<html lang>}, \texttt{<head>} avec \texttt{charset} et \texttt{<title>}, \texttt{<body>}
      \item Un \texttt{<header>} avec un \texttt{<h1>} portant votre prénom
      \item Un \texttt{<main>} avec une \texttt{<section>} : un \texttt{<h2>}, un paragraphe, un mot en \texttt{<strong>}
      \item Une image, avec un \texttt{alt} correct
      \item Un lien vers un site que vous aimez, à l'intitulé explicite
      \item Un \texttt{<footer>} avec la date du jour
    \end{enumerate}
  \end{block}
  \begin{center}
    \small\itshape 20 minutes --- puis on valide ensemble sur \texttt{validator.w3.org}.
  \end{center}
\end{frame}

\begin{frame}{Exercice --- les points d'attention}
  \begin{columns}[T]
    \begin{column}{0.49\textwidth}
      \begin{exampleblock}{Si vous êtes en avance}
        \begin{itemize}
          \item Ajouter une liste \texttt{<ul>} de trois centres d'intérêt
          \item Ajouter un second \texttt{<h2>} et une deuxième \texttt{<section>}
          \item Placer l'image dans une \texttt{<figure>} avec sa \texttt{<figcaption>}
        \end{itemize}
      \end{exampleblock}
    \end{column}
    \begin{column}{0.49\textwidth}
      \begin{alertblock}{Ce que je vérifierai en passant}
        \begin{itemize}
          \item Le code est \textbf{indenté}
          \item \textbf{Un seul} \texttt{<h1>}
          \item L'image a un \texttt{alt} qui décrit quelque chose
          \item Le chemin de l'image est \textbf{relatif}
          \item Aucune balise n'est laissée ouverte
        \end{itemize}
      \end{alertblock}
    \end{column}
  \end{columns}
  \vspace{0.25cm}
  \begin{center}
    \small \textbf{Pas de squelette généré automatiquement} pour cet exercice : on l'écrit à la main, une fois.\\
    \small Le raccourci \texttt{!} + \texttt{Tab} de VS Code sera autorisé dès le TP.
  \end{center}
\end{frame}

\begin{frame}[fragile]{Exercice --- une correction possible}
\begin{lstlisting}[language=HTML5]
<!DOCTYPE html>
<html lang="fr">
<head>
  <meta charset="UTF-8">
  <title>Alice - présentation</title>
</head>
<body>
  <header><h1>Alice</h1></header>
  <main>
    <section>
      <h2>Qui suis-je ?</h2>
      <p>Étudiante en classe passerelle, je découvre le
         <strong>développement web</strong> cette année.</p>
      <img src="images/photo.jpg" alt="Alice devant l'entrée du cfpt">
      <p><a href="https://developer.mozilla.org">La documentation MDN</a></p>
    </section>
  </main>
  <footer><p>Page créée le 15 septembre 2026</p></footer>
</body>
</html>
\end{lstlisting}
\end{frame}

% ===============================================================
\section{Listes \& tableaux}
% ===============================================================
\partieslide{6. Listes \& tableaux}

\begin{frame}[fragile]{Les listes}
  \begin{columns}[T]
    \begin{column}{0.5\textwidth}
      \begin{center}\footnotesize\textbf{Liste à puces --- ordre indifférent}\end{center}
\begin{lstlisting}[language=HTML5]
<ul>
  <li>HTML</li>
  <li>CSS</li>
  <li>JavaScript</li>
</ul>
\end{lstlisting}
    \end{column}
    \begin{column}{0.5\textwidth}
      \begin{center}\footnotesize\textbf{Liste numérotée --- l'ordre compte}\end{center}
\begin{lstlisting}[language=HTML5]
<ol>
  <li>Ouvrir le terminal</li>
  <li>Taper la commande</li>
  <li>Valider</li>
</ol>
\end{lstlisting}
    \end{column}
  \end{columns}
  \vspace{0.25cm}
  \begin{alertblock}{Un \texttt{<li>} vit toujours dans un \texttt{<ul>} ou un \texttt{<ol>}}
    Rien d'autre qu'un \texttt{<li>} ne peut être enfant direct d'un \texttt{<ul>} : pas de \texttt{<p>}, pas de \texttt{<div>}. Ce contenu se place \emph{à l'intérieur} du \texttt{<li>}.
  \end{alertblock}
\end{frame}

\begin{frame}[fragile]{Un menu de navigation est une liste de liens}
\begin{lstlisting}[language=HTML5]
<nav>
  <ul>
    <li><a href="index.html">Accueil</a></li>
    <li><a href="sujet.html">Sujets</a></li>
  </ul>
</nav>
\end{lstlisting}
  \vspace{0.3cm}
  \begin{block}{Étrange ? Pas tant que ça}
    Cela paraît lourd tant qu'on pense en termes d'\textbf{apparence}. Mais un menu \emph{est} bien une liste de liens --- et le CSS se chargera, en semaine 4, de l'afficher horizontalement.
  \end{block}
  \vspace{0.2cm}
  \begin{center}
    \small C'est exactement l'esprit du HTML : \textbf{décrire le sens}, laisser le CSS décider de la forme.
  \end{center}
\end{frame}

\begin{frame}[fragile]{Les tableaux}
\begin{lstlisting}[language=HTML5]
<table>
  <caption>Derniers sujets publiés</caption>
  <thead>
    <tr><th>Sujet</th><th>Auteur</th></tr>
  </thead>
  <tbody>
    <tr><td>Bienvenue</td><td>Alice</td></tr>
    <tr><td>Installer Node.js</td><td>Bob</td></tr>
  </tbody>
</table>
\end{lstlisting}
  \vspace{0.15cm}
  \begin{alertblock}{Un tableau sert à des \textbf{données}, jamais à mettre en page}
    Horaire, relevé de notes, comparatif : oui. Positionner des éléments dans la page : non --- c'est le rôle du CSS. Une page mise en page avec des tableaux est illisible pour un lecteur d'écran.
  \end{alertblock}
\end{frame}

% ===============================================================
\section{Valider son HTML}
% ===============================================================
\partieslide{7. Valider son HTML}

\begin{frame}{Pourquoi valider}
  \begin{block}{Le navigateur est trop gentil}
    Il affiche une page même si des balises sont mal fermées : il \textbf{devine} ce que vous avez voulu écrire.
  \end{block}
  \vspace{0.3cm}
  \begin{alertblock}{Et c'est un piège}
    Chaque navigateur devine à sa manière. Le jour où l'affichage se dérègle, l'erreur est déjà enfouie sous cent lignes de code --- et elle ne se manifeste que chez certains utilisateurs.
  \end{alertblock}
  \vspace{0.3cm}
  \begin{exampleblock}{Le validateur du W3C --- \texttt{validator.w3.org}}
    Trois modes : par \textbf{adresse}, par \textbf{fichier} téléversé, ou par \textbf{copier-coller} --- c'est celui qu'on utilise en classe.
  \end{exampleblock}
\end{frame}

\begin{frame}{Lire un message d'erreur}
  Chaque message comporte trois parties : la \textbf{ligne et la colonne}, la \textbf{description}, et un \textbf{extrait} du code fautif.
  \vspace{0.25cm}
  \begin{center}
  \begin{tikzpicture}
    \node[draw=cfptDarkGrey!50, fill=cfptGrey, rounded corners=2pt,
          inner sep=8pt, text width=0.86\textwidth, font=\ttfamily\small] {
      Error: End tag \textcolor{alertRed}{section} seen, but there were open elements.\\
      From line 24, column 3; to line 24, column 12};
  \end{tikzpicture}
  \end{center}
  \vspace{0.25cm}
  \begin{block}{Traduction}
    À la ligne 24, on ferme une \texttt{<section>} alors qu'un élément ouvert à l'intérieur ne l'a pas été. Il manque une balise fermante \textbf{au-dessus}.
  \end{block}
\end{frame}

\begin{frame}{Les messages les plus fréquents}
  \begin{center}
  \small
  \begin{tabularx}{0.97\textwidth}{p{5.6cm}X}
    \toprule
    \rowcolor{accentBlue}
    \hdr{Message} & \hdr{Ce qu'il faut corriger} \\
    \midrule
    \rowcolor{cfptGrey}
    \texttt{End tag X seen, but there were open elements} & Une balise ouverte à l'intérieur n'a pas été fermée \\
    \texttt{Unclosed element X} & La balise \texttt{X} n'est jamais fermée \\
    \rowcolor{cfptGrey}
    \texttt{Duplicate ID "xyz"} & Deux éléments portent le même \texttt{id} \\
    \texttt{An img element must have an alt attribute} & Une image sans \texttt{alt} \\
    \rowcolor{cfptGrey}
    \texttt{Element X not allowed as child of element Y} & Mauvaise imbrication \\
    \texttt{Stray end tag X} & Une fermante sans ouvrante \\
    \bottomrule
  \end{tabularx}
  \end{center}
  \vspace{0.25cm}
  \begin{alertblock}{Méthode --- corriger de haut en bas}
    Une seule balise oubliée peut produire dix messages en cascade. Corrigez \textbf{la première erreur}, relancez : le compteur chute souvent d'un coup. Les \textbf{Errors} doivent toutes disparaître ; les \textbf{Warnings} sont des maladresses.
  \end{alertblock}
\end{frame}

\begin{frame}{Prévenir plutôt que corriger}
  \begin{block}{Trois habitudes qui suppriment la majorité des erreurs}
    \begin{itemize}
      \item \textbf{Indenter} systématiquement --- \texttt{Maj+Alt+F} dans VS Code, avec Prettier
      \item Écrire la balise \textbf{fermante immédiatement} après l'ouvrante, puis remplir le contenu
      \item Se fier au \textbf{soulignement rouge} de VS Code, qui repère déjà une partie des erreurs
    \end{itemize}
  \end{block}
  \vspace{0.4cm}
  \begin{exampleblock}{Le réflexe documentation}
    \textbf{MDN Web Docs} reste la première source : cherchez \texttt{mdn} suivi du nom de la balise. La page donne la description, la syntaxe, et surtout des \textbf{exemples modifiables}.
  \end{exampleblock}
\end{frame}

% ===============================================================
% Clôture
% ===============================================================
\begin{frame}{Résumé de la semaine}
  \begin{columns}[T]
    \begin{column}{0.5\textwidth}
      \begin{block}{Ce qu'on a vu}
        \begin{itemize}
          \item Le HTML, langage de balisage
          \item L'imbrication, et l'arbre du document
          \item Le squelette : \texttt{<head>} / \texttt{<body>}
          \item Les balises sémantiques
          \item Titres, texte, liens, images
          \item Listes et tableaux
          \item La validation W3C
        \end{itemize}
      \end{block}
    \end{column}
    \begin{column}{0.5\textwidth}
      \begin{alertblock}{À faire pour la semaine 3}
        \begin{itemize}
          \item Les \textbf{deux pages} HTML du MiniForum (fiche TP)
          \item Chacune valide \textbf{sans erreur}
          \item Pousser le tout sur GitHub
        \end{itemize}
      \end{alertblock}
      \vspace{0.15cm}
      \begin{exampleblock}{La semaine prochaine}
        Les \textbf{formulaires} et l'\textbf{accessibilité} --- encore du HTML, et la dernière semaine avant le CSS.
      \end{exampleblock}
    \end{column}
  \end{columns}
\end{frame}

{
\setbeamertemplate{footline}{}
\begin{frame}[plain]
  \begin{tikzpicture}[remember picture, overlay]
    \fill[accentBlue] (current page.north west) rectangle (current page.south east);
    \node[text=white, font=\Huge\bfseries] at ([yshift=0.6cm]current page.center) {Des questions ?};
    \fill[cfptYellow] ([yshift=-0.2cm]current page.center -| current page.west)
      rectangle ([yshift=-0.32cm]current page.center -| current page.east);
    \node[text=lightBlue, font=\large] at ([yshift=-1.3cm]current page.center)
      {Prochaine étape : la fiche TP --- structure HTML du MiniForum};
  \end{tikzpicture}
\end{frame}
}

\end{document}
